\section{Introduction}
\label{sec:introduction}

Recent advances in video generation~\cite{ho2020denoising,kong2024hunyuanvideo,wan2025wan,polyak2024moviegen,brooks2024sora} have enabled interactive world models with controllable generation across games~\cite{bruce2024genie,deepmind2024genie2,ball2025genie3,he2025matrixgame2,he2026matrixgame3,li2025hunyuangamecraft,tang2025interbench}, autonomous driving~\cite{hu2023gaia1,gao2024vista}, embodied interaction~\cite{yang2023unisim,zhu2024irasim}, and open-domain scenarios~\cite{sun2025hyworldplay,chen2025yume,mao2025yume15,gao2026lingbot}.
However, evaluation remains fragmented, with many works relying on selected demos or task-specific protocols, making fair comparison and failure diagnosis difficult across visual quality, controllability, memory, and physics.

A capable interactive world model must fulfill five complementary roles, analogous to the subsystems of a game engine: a \emph{Renderer} for visually convincing video, a \emph{Director} for correct world initialization, a \emph{Controller} for faithful interaction execution, a \emph{Memory} for preserving world state across turns, and an \emph{Engine} for physically compliant world evolution. Existing benchmarks cover these roles only partially (\cref{tab:benchmark_comparison}). Video-generation benchmarks such as VBench~\cite{huang2024vbench,zheng2025vbench2} focus on perceptual quality without interactive control. World-model benchmarks evaluate more dimensions but remain limited in scope: WorldMark~\cite{xu2026worldmark} and MIND~\cite{ye2026mind} cover navigation and memory but lack semantic interactions, Omni-WorldBench~\cite{wu2026omniworldbench} adds causal interaction but supports only first-person view, and WorldLens~\cite{worldlens2025} evaluates multiple dimensions but is restricted to autonomous driving. None provides a unified protocol spanning open-domain scenes, both perspectives, and all four interaction types.

To address this gap, we introduce \benchmark, a comprehensive multi-turn benchmark for interactive world model evaluation. As shown in \cref{fig:teaser}, each test case is defined by a \emph{world setting} (scene, subject, style, and perspective) together with a multi-turn \emph{interaction} sequence. The top row illustrates a concrete case: a realistic snowy mountain scene with a human subject in third-person perspective, followed by forward navigation, a jump, the appearance of a helicopter, and a perspective switching to the cockpit. More broadly, the benchmark spans diverse open-domain scenes, rendering styles, subject categories, and both first- and third-person perspectives (\cref{fig:teaser}~\!(a)), with four interaction types shown in \cref{fig:teaser}~\!(b): navigation, subject action, event editing, and perspective switching. This design separates what the world \emph{is} from what the user \emph{requests}, making failure modes easier to locate: a model may render the initial scene well but ignore later actions, or follow a single instruction correctly but lose identity and spatial consistency over multiple turns.

\benchmark also supports fair comparison across different control paradigms. As shown in \cref{fig:teaser}~\!(c), navigation interactions are represented in three aligned forms, namely text, camera pose, and discrete action, so that models can be evaluated through their native interfaces. Accordingly, we adopt a dual-track evaluation protocol: all \nummodel models are compared on a shared navigation subset of \numnavi cases, while text-prompted I2V models are further evaluated on the full benchmark (\numvideo cases, \numturn turns). Evaluation uses \numsubmetric automatic sub-metrics combining specialist vision models and VLMs.

Experiments on \nummodel models reveal that:
1)~no model dominates all five dimensions,
2)~navigation is largely independent of other dimensions,
3)~camera control and perspective consistency are separate capabilities,
4)~physical correctness correlates with rendering quality rather than control,
5)~benchmark difficulty is structured by perspective, scene type, and subject category, and
6)~four interaction types degrade unevenly over turns, with navigation most fragile.

Our contributions are: 1)~a unified benchmark spanning five complementary evaluation dimensions with \numsubmetric fine-grained sub-metrics, 2)~a multi-turn dataset covering both perspectives, four interaction types, and a unified navigation interface enabling fair cross-paradigm comparison, and 3)~a fully automatic evaluation pipeline applied to \nummodel models, establishing diagnostic baselines and surfacing actionable insights for future model development.
